\corrigesChapitre{14}{Mesurer l’intensité du courant}
\par\noindent\begin{minipage}{\linewidth}
\begin{corrige}{1}{L’unité}
\begin{questions}
\item L’ampère, A.
\item L’ampèremètre, généralement une fonction du multimètre.
\end{questions}
\end{corrige}
\end{minipage}
\par\noindent\begin{minipage}{\linewidth}
\begin{corrige}{2}{Le branchement}
\begin{questions}
\item En série dans la branche étudiée.
\item Oui, puis faire vérifier le montage selon les consignes.
\end{questions}
\pcFigureEx{figures/cor-amperemetre.pdf}{L’ampèremètre est inséré dans la branche étudiée ; couper l’alimentation avant de modifier le montage.}
\end{corrige}
\end{minipage}
\par\noindent\begin{minipage}{\linewidth}
\begin{corrige}{3}{Convertir}
\begin{questions}
\item $0{,}250\unite A$.
\item $80\unite{mA}$.
\end{questions}
\end{corrige}
\end{minipage}
\par\noindent\begin{minipage}{\linewidth}
\begin{corrige}{4}{Dans une série}
\begin{questions}
\item $0{,}12\unite A$ dans le modèle en régime établi.
\item Non ; elle a converti de l’énergie.
\end{questions}
\pcFigureEx{figures/cor-serie-012.pdf}{En régime permanent, même intensité avant et après la lampe dans cette unique branche.}
\end{corrige}
\end{minipage}
\par\noindent\begin{minipage}{\linewidth}
\begin{corrige}{5}{Un nœud simple}
\begin{questions}
\item $0{,}30\unite A$.
\item L’additivité des intensités au nœud, avec les sens annoncés.
\end{questions}
\pcFigureEx{figures/cor-noeud.pdf}{La somme des courants entrants est égale à la somme des courants sortants.}
\end{corrige}
\end{minipage}
\par\noindent\begin{minipage}{\linewidth}
\begin{corrige}{6}{Une branche inconnue}
\begin{questions}
\item $I_2=0{,}50-0{,}18=0{,}32\unite A$.
\item $0{,}18+0{,}32=0{,}50\unite A$.
\end{questions}
\par\smallskip\noindent\textbf{Pour comprendre.}\enspace Au nœud, $I_G=I_1+I_2$. Isoler $I_2=I_G-I_1=0{,}50-0{,}18=0{,}32\unite A$. Contrôle : $0{,}18+0{,}32=0{,}50\unite A$.\par
\pcFigureEx{figures/cor-courants-050.pdf}{Les intensités sont exprimées dans la même unité avant la soustraction.}
\end{corrige}
\end{minipage}
\par\noindent\begin{minipage}{\linewidth}
\begin{corrige}{7}{Attention aux unités}
\begin{questions}
\item $I_1=250\unite{mA}$ ; $I_2=600-250=350\unite{mA}$.
\item $0{,}350\unite A$.
\end{questions}
\par\smallskip\noindent\textbf{Pour comprendre.}\enspace Employer une seule unité : $0{,}25\unite A=250\unite{mA}$. Puis $I_2=600-250=350\unite{mA}=0{,}350\unite A$.\par
\end{corrige}
\end{minipage}
\par\noindent\begin{minipage}{\linewidth}
\begin{corrige}{8}{Choisir un calibre}
\begin{questions}
\item 200 mA est adapté si les bornes, protections et conditions d’utilisation le permettent.
\item La valeur réelle peut être plus élevée ; on commence selon la notice et les consignes, avec un calibre sûr.
\end{questions}
\end{corrige}
\end{minipage}
\par\noindent\begin{minipage}{\linewidth}
\begin{corrige}{9}{Le signe négatif}
\begin{questions}
\item Il circule dans le sens opposé à la convention positive du branchement.
\item Sa valeur est $0{,}15\unite A$ ; on vérifie les connexions hors alimentation.
\end{questions}
\pcFigureEx{figures/cor-amperemetre-inverse.pdf}{Le signe négatif provient de l’orientation de connexion retenue ; la valeur absolue est 0,15 A.}
\end{corrige}
\end{minipage}
\par\noindent\begin{minipage}{\linewidth}
\begin{corrige}{10}{Deux branches identiques}
\begin{questions}
\item $0{,}20\unite A$ par branche.
\item Non : elle suppose des branches identiques et les mêmes conditions.
\end{questions}
\end{corrige}
\end{minipage}
\par\noindent\begin{minipage}{\linewidth}
\begin{corrige}{11}{Une expérience cohérente ?}
\begin{questions}
\item $0{,}122+0{,}176=0{,}298\unite A$ ; l’écart vaut $0{,}002\unite A$.
\item Non : comparer à la précision, aux variations de la source et aux conditions de mesure.
\end{questions}
\par\smallskip\noindent\textbf{Pour comprendre.}\enspace La somme mesurée vaut $0{,}122+0{,}176=0{,}298\unite A$ ; l’écart à 0,300 A est 0,002 A, soit 2 mA. Comparer cet écart à la précision et à la stabilité des mesures avant de conclure.\par
\end{corrige}
\end{minipage}
\par\noindent\begin{minipage}{\linewidth}
\begin{corrige}{12}{Un emplacement dangereux}
\begin{questions}
\item L’appareil de faible résistance peut court-circuiter la pile.
\item Il faut l’insérer dans la branche dont on veut mesurer le courant, hors alimentation et avec un calibre adapté.
\end{questions}
\pcFigureEx{figures/cor-amperemetre.pdf}{L’ampèremètre doit mesurer le courant d’une branche contenant le récepteur, pas relier seul les bornes de la pile.}
\end{corrige}
\end{minipage}
\par\noindent\begin{minipage}{\linewidth}
\begin{corrige}{13}{Une lampe supplémentaire}
\begin{questions}
\item Non : le circuit a été modifié et sa valeur de courant peut changer.
\item La loi compare les intensités en différents points d’un même circuit en régime établi.
\end{questions}
\end{corrige}
\end{minipage}
\par\noindent\begin{minipage}{\linewidth}
\begin{corrige}{14}{Une branche ouverte}
\begin{questions}
\item Le courant de la branche 2 devient nul et $I_G=I_1$ pour la branche encore fermée.
\item Non : cela dépend aussi du générateur ; le modèle idéal le permet, mais une source réelle peut varier.
\end{questions}
\pcFigureEx{figures/cor-courant-branche-ouverte.pdf}{La branche ouverte a une intensité nulle ; le courant principal est celui de la seule branche encore fermée.}
\end{corrige}
\end{minipage}
\par\noindent\begin{minipage}{\linewidth}
\begin{corrige}{15}{Trois branches}
\begin{questions}
\item $I_3=0{,}80-0{,}220-0{,}35=0{,}230\unite A$.
\item Les trois valeurs sont positives et leur somme vaut $0{,}80\unite A$.
\end{questions}
\par\smallskip\noindent\textbf{Pour comprendre.}\enspace Convertir $220\unite{mA}=0{,}220\unite A$. $I_3=0{,}80-0{,}220-0{,}35=0{,}23\unite A$ ; contrôle : $0{,}220+0{,}35+0{,}23=0{,}80\unite A$.\par
\pcFigureEx{figures/cor-trois-courants.pdf}{Convertir 220 mA en 0,220 A ; soustraire les deux branches connues du courant principal.}
\end{corrige}
\end{minipage}
\par\noindent\begin{minipage}{\linewidth}
\phantomsection\label{cor-defi-14}\pcSolutionDefi{14}{Trouver la mesure incohérente}\begin{questions}
\item $0{,}260+0{,}24=0{,}50\unite A$, ce qui diffère de $0{,}40\unite A$.
\item $0{,}50\unite A$.
\item Vérifier les unités et calibres, le placement des appareils, et la stabilité du générateur entre les mesures ; contrôler aussi l’existence d’autres branches.
\end{questions}
\pcFigureEx{figures/cor-bilan-intensites.pdf}{Le schéma relie chaque valeur à sa branche ; les mesures fournies ne respectent pas l’additivité dans leurs conditions affichées.}
\end{minipage}
\par\noindent\begin{minipage}{\linewidth}
\phantomsection\label{cor-activite-14}\pcSousTitre{Pistes pour l’investigation}
Dans le montage série, les mesures avant et après la lampe sont compatibles avec la même intensité dans la limite des appareils. En dérivation, comparer le courant principal à la somme des courants des branches. Un écart nécessite une discussion des calibres, des connexions et de la stabilité de la source, sans déclarer automatiquement une loi fausse.\par\end{minipage}
